\lyseisE{Επανάληψη κεφαλαίου Β1}
\solr{Β1}{1} Τρεις ευθείες και έξι ευθείες.
\solr{Β1}{2} Πεντάγωνο $5$, εξάγωνο $9$.
\solr{Β1}{3} Το μέσο απέχει $3{,}6$ cm από κάθε άκρο. Επειδή $7{,}2$ cm $=72$ mm $>70$ mm, το $\mathrm{ΑΒ}$ είναι μεγαλύτερο.
\solr{Β1}{4} $\mathrm{ΑΔ}=2{,}5+3+4=9{,}5$ και $\mathrm{ΒΔ}=3+4=7$.
\solr{Β1}{5} $2(6{,}5+4)=21$.
\solr{Β1}{6} Διχοτόμος: κάθε μισή $62^\circ$. Παραπλήρωμα $56^\circ$. Συμπλήρωμα της $36^\circ$: $54^\circ$.
\solr{Β1}{7} $65^\circ$, $115^\circ$ και $115^\circ$.
\solr{Β1}{8} $360^\circ-(90+80+110)^\circ=80^\circ$.
\solr{Β1}{9} $150^\circ$, αμβλεία.
\solr{Β1}{10} Το πολύ $\frac{4\cdot3}{2}=6$ σημεία. Μία παράλληλη.
\solr{Β1}{11} $1{,}8+1{,}8=3{,}6$ cm.
\solr{Β1}{12} Διάμετρος $12$. Το $\mathrm{Α}$ μέσα ($4<6$), το $\mathrm{Β}$ έξω ($7>6$).
\solr{Β1}{13} $360:9=40^\circ$. Το τόξο είναι $150^\circ$.
\solr{Β1}{14} α) εξωτερική ($0$ κοινά) \quad β) εφαπτομένη ($1$) \quad γ) τέμνουσα ($2$).
\solr{Β1}{15} α) Σωστό. β) Λάθος. γ) Σωστό. δ) Σωστό. ε) Σωστό.
\solr{Β1}{16} α) Σωστό. β) Λάθος: $90^\circ$. γ) Σωστό. δ) Σωστό. ε) Σωστό. στ) Λάθος: είναι το μικρότερο.
